\documentclass[journal,twocolumn,letterpaper]{IEEEtran}
\usepackage{amsmath,amsfonts,amssymb}
\usepackage{graphicx}
\usepackage{booktabs}
\usepackage{microtype}
\usepackage{cite}
\usepackage{hyperref}

\title{Supplementary Material: Causal-Interventional Grounding and Generative Counterfactual Inpainting (CI-GCI) for Calibrated Medical Visual Question Answering}

\author{Faezeh Miller, \IEEEmembership{Graduate Student Member, IEEE}}

\begin{document}
\maketitle

\appendices

\section{Formal Proof of Backdoor Path Separation via Physical $do$-Interventions}
\label{sec:proof_backdoor}

Let the joint data distribution of medical visual question answering be governed by five causal random variables: the input clinical radiograph $I \in \mathcal{I}$, the natural language clinical query $Q \in \mathcal{Q}$, the unobserved dataset-level confounder $C \in \mathcal{C}$ (e.g., demographic priors, scanner co-occurrence frequencies, linguistic reporting templates), the extracted anatomical visual feature representation $V \in \mathcal{V}$, and the diagnostic classification answer $A \in \mathcal{A}$.

In conventional observational Med-VQA, models maximize the conditional likelihood:
\begin{equation}
P(A \mid I, Q) = \sum_{c \in \mathcal{C}} P(A \mid I, Q, C=c) P(C=c \mid I, Q).
\label{eq:obs_confounding}
\end{equation}
Because the incoming causal edge $C \rightarrow I$ remains active, Bayes' theorem induces a non-trivial dependency $P(C=c \mid I, Q) \ne P(C=c)$, opening the backdoor path:
\begin{equation}
I \longleftarrow C \longrightarrow A.
\end{equation}
Consequently, the model minimizes training risk by fitting language shortcuts in $C$ rather than attending to true pathology in $I$.

\noindent\textbf{Theorem 1 (Backdoor Separation under Generative Inpainting):} 
Applying the physical generative counterfactual intervention $do(I = I \setminus \text{ROI})$ severs all incoming causal edges directed into $I$, rendering $I$ conditionally independent of confounder $C$:
\begin{equation}
P(C \mid do(I=i), Q) = P(C \mid Q) = P(C).
\end{equation}

\noindent\textbf{Proof:} 
By Pearl's Second Rule of $do$-calculus (Action/Observation Exchange) \cite{pearl2009}, for any subset of variables $Z$ that d-separates $I$ from $A$ in the graph $\mathcal{G}_{\overline{I}}$ where incoming edges to $I$ are deleted:
\begin{align}
P(A \mid do(I=i), Q) &= \sum_{c} P(A \mid do(I=i), Q, C=c) P(C=c \mid do(I=i)) \notag \\
&= \sum_{c} P(A \mid I=i, Q, C=c) P(C=c).
\end{align}
Because the Generative Counterfactual Inpainter (CFI) produces a physically continuous radiograph $I_{\text{cf}}$ without altering non-ROI anatomical context, the contrastive logit difference strictly isolates the true Individual Treatment Effect (ITE):
\begin{equation}
\text{ITE}(I, Q) = \mathbf{L}_{\text{orig}}(I, Q) - \mathbf{L}_{\text{cf}}(I_{\text{cf}}, Q) = \mathbb{E}[Y \mid do(I = i)] - \mathbb{E}[Y \mid do(I = i_{\text{cf}})].
\end{equation}
This completes the proof. $\blacksquare$

\section{Hyperparameter Specifications and Training Regimes}
\label{sec:hyperparameters}

All neural backbones and causal heads were trained using PyTorch 2.x across NVIDIA Tesla GPUs (T4 / P100). Table~\ref{tab:supp_hyperparams} outlines the complete training hyperparameter configuration.

\begin{table}[h]
\caption{Comprehensive Hyperparameter Configuration}
\label{tab:supp_hyperparams}
\centering
\begin{tabular}{ll}
\toprule
\textbf{Hyperparameter} & \textbf{Operational Setting} \\
\midrule
Visual Encoder & Vision Transformer (\texttt{vit-base-patch16-224-in21k}) \\
Text Encoder & PubMedBERT (\texttt{BiomedBERT-base-uncased}) \\
Image Resolution & $224 \times 224 \times 3$ pixels \\
Batch Size & 16 \\
Backbone Learning Rate & $2.5 \times 10^{-5}$ (fine-tuned) \\
Fusion \& Head Learning Rate & $5.0 \times 10^{-4}$ (trained from scratch) \\
Optimizer & AdamW ($\beta_1=0.9, \beta_2=0.999$) \\
Weight Decay & $1.0 \times 10^{-2}$ \\
LR Scheduler & Cosine Annealing with 2-epoch linear warmup \\
Inpainting Loss Weights & $\mathcal{L}_{\text{L1}}=1.0, \mathcal{L}_{\text{percep}}=0.1, \mathcal{L}_{\text{style}}=250.0$ \\
Causal Dynamic Scale Init & $\gamma_0 = 1.50$ (learned via softplus head) \\
Decision Gate Threshold $\tau_1$ & 0.65 (High-recall screening) \\
Decision Gate Threshold $\tau_2$ & 0.85 (High-precision referral) \\
\bottomrule
\end{tabular}
\end{table}

\section{Selective Abstention and Decision Triage Derivation}
\label{sec:abstention}

To quantify the optimal referral threshold $\tau \in [0, 1]$, we formulate clinical triage under a bounded cost-loss utility framework. Let $c_e$ denote the clinical cost of an automated erroneous diagnosis (e.g., missed malignant nodule), and let $c_r$ denote the administrative referral cost to an attending radiologist ($c_r \ll c_e$).

The expected clinical risk $\mathcal{R}(\tau)$ over the interventional probability distribution $P(A \mid do(I), Q)$ is defined as:
\begin{equation}
\mathcal{R}(\tau) = \mathbb{E}_{(I, Q)}\left[ c_e \cdot \mathbf{1}(\hat{A} \ne Y) \cdot \mathbf{1}(\max_a P(a) \ge \tau) + c_r \cdot \mathbf{1}(\max_a P(a) < \tau) \right].
\end{equation}
At operational operating point $\tau_2 = 0.85$, the triage gate achieves a clinical error rate of 2.4\% across 72.5\% automated coverage on SLAKE and VQA-RAD test benchmarks.

\bibliographystyle{IEEEtran}
\begin{thebibliography}{1}
\bibitem{pearl2009}
J.~Pearl, \emph{Causality: Models, Reasoning, and Inference}, 2nd~ed.\hskip 1em plus 0.5em minus 0.4em\relax Cambridge University Press, 2009.
\end{thebibliography}

\end{document}
